% Tabela (agente kl_talude): custo das formulacoes no talude. Fontes: Galerkin, dados/kl_galerkin_talude_regra_elemento.csv e
% dados/kl_galerkin_talude_ordem10.csv (rodadas refeitas para o relatorio, KLTaludeCeron, 1 thread);
% Nystrom, dados/kl_nystrom_custo.csv; forma fechada, README do KLCilindroMika, secao 4.
\begin{table}[H]
  \centering
  \caption{Custo de cada formulação no talude (uma thread). Galerkin: KL numérica do NeoPZ (\texttt{KLTaludeCeron}), tempo
    de montagem de $\mathbf{C}$, do problema de autovalores com autovetores e total, e memória máxima, nas rodadas refeitas
    para este relatório; os tempos dependem da carga da máquina (o README registra 260~s e 650~MB para a malha de 5085
    equações com a regra de ordem 10).
    Nyström: $n = 160$, 210 autopares. Forma fechada: 20\,000 termos e as integrais de cada modo no talude.}
  \label{tab:kl:custo}
  \footnotesize
  \setlength{\tabcolsep}{4pt}
  \begin{tabular}{llcrrrrr}
    \toprule
    KL & malha & $P$ / regra & incógnitas & $\mathbf{C}$ (s) & autopares (s) & total (s) & memória (MB) \\
    \midrule
    Galerkin & \texttt{ceron\_h5\_a45}      & 1 / 2  &    609 &  0,11 &   0,22 &   0,71 &  61 \\
    Galerkin & \texttt{ceron\_h5\_a45\_h025} & 1 / 2  &   1314 &  0,58 &   1,99 &   3,36 &  91 \\
    Galerkin & \texttt{ceron\_h5\_a45\_fina} & 1 / 2  &   2301 &  2,39 &  12,03 &  16,10 & 206 \\
    Galerkin & \texttt{ceron\_h5\_a45}      & 2 / 4  &   2320 &  0,46 &  12,24 &  13,33 & 209 \\
    Galerkin & \texttt{ceron\_h5\_a45\_h025} & 2 / 4  &   5085 &  2,96 & 199,83 & 204,06 & 646 \\
    Galerkin & \texttt{ceron\_h5\_a45\_fina} & 1 / 10 &   2301 & 62,15 &  10,84 &  76,67 & 204 \\
    Galerkin & \texttt{ceron\_h5\_a45}      & 2 / 10 &   2320 &  4,43 &  12,10 &  17,55 & 204 \\
    Galerkin & \texttt{ceron\_h5\_a45\_h025} & 2 / 10 &   5085 & 22,44 & 190,99 & 215,57 & 653 \\
    \midrule
    Nyström  & grade $h = \SI{3.125}{cm}$    & --     & 193\,041 & --  &  77,3  &  80,9  & 1414 \\
    forma fechada & --                       & --     & --     & --    & --     &   1,3  & -- \\
    \bottomrule
  \end{tabular}
\end{table}
